\section{反標準環による包絡の一意性とAS曲面との比較}
\label{sec:as-surface-scope}

本節では、四次元AS正則環の三次商であることと、
三次元AS正則$\mathbb Z$-代数の$\qgr$であることを区別する。
後者がOkawa--Ueda \cite[Definition~1.1]{OUASSurfaces}のAS曲面であり、
del Pezzo型の標識には対応するvery strong helixが必要である。
以下の一意性は、構造対象を保つexact同値に対する主張である。

\begin{lemma}[飽和反標準環と二次包絡]\label{lem:anticanonical-quadratic-hull}
構造対象を持つ二つの圏$(\mathcal A,\OO)$、$(\mathcal A',\OO')$が
Serre作用を持ち、$\tau=S^{-1}[2]$、$\tau'=(S')^{-1}[2]$とする。
構造対象を保つexact同値は、反標準orbit環
\[
 R=\bigoplus_{n\geq0}\Hom(\OO,\tau^n\OO),\qquad
 R'=\bigoplus_{n\geq0}\Hom(\OO',(\tau')^n\OO')
\]
の次数付き同型を誘導する。
さらに$R=B/(G)$で、$B$が一次生成二次関係環、$G$が三次斉次元なら、
\[
 Q(R):=T_k(R_1)/\bigl(\ker(R_1\otimes_kR_1\longrightarrow R_2)\bigr)
 \simeq B.
\]
従って、このような二次包絡の次数付き同型類は$(\mathcal A,\OO)$から決まる。
\end{lemma}
\begin{proof}
exact同値の導来拡張はSerre作用をその一意性により保つ。
構造対象の同型と$F\tau\simeq\tau'F$の反復を用いると、
各次数のHomの同型がorbit積と両立する。
後半では、三次元による商は次数一・二の部分とその乗法を変えない。
よって$R_1=B_1$であり、二次乗法の核は$B$の二次関係空間そのものである。
$B$はこの関係空間で定義されるので結論を得る。
\end{proof}

この補題を斉次表示へ適用するには、表示した商環が
\emph{飽和反標準環そのものであること}を確認する必要がある。
任意の次数付き環の$\qgr$同値、任意の導来同値、
構造対象を指定しない同値からの包絡復元は、ここでは主張しない。

\begin{proposition}[AOUの無限位数blow-up表示との分離]
\label{prop:aou-infinite-order-exclusion}
式\eqref{eq:B}の$B_{\mathbf a}$と、そのH3切断
$R_{\mathbf a,u}=B_{\mathbf a}/(G_{\mathbf a,u})$を取る。
入力\ref{ass:recognition}(4)に従い、この商環が対応するpointed圏の
飽和反標準環であるとする。
このpointed圏は、Abdelgadir--Okawa--Ueda
\cite[\S9]{AOUCubic}で\emph{無限位数の平行移動}から構成される
反標準商$(\qgr F,\pi F)$とpointed exact同値ではない。
\end{proposition}
\begin{proof}
同文献の式(9.6)--(9.10)は$F=P/(C'_3)$という二次AS4包絡を与え、
Lemma~9.2は逆Serre作用を$(1)$と同定する。
AS Cohen--Macaulay性による飽和性と合わせて$Q(F)=P$を得る。
同節では中心一次直線は$kt$であり、$P/(t)$は滑らかな楕円曲線を
point schemeとする三次元Sklyanin環である。

一方、命題\ref{prop:pbw}より$Z(B_{\mathbf a})_1=kt$で、
\[
 B_{\mathbf a}/(t)=S_q
 =k\langle x,y,z\rangle/(yz-qzy,\ zx-qxz,\ xy-qyx).
\]
$q^3=1$かつ$q\ne1$なので、$S_q$の二次関係を多重線型化したschemeは
$\mathbb P^2$全体上の$\phi=\operatorname{diag}(1,q,q^2)$のgraphである。
実際、$(x',y',z')=(x,qy,q^2z)$は全関係を満たし、
任意の非零$(x,y,z)$でこの三線型方程式の係数行列はrank~2である。
したがってその核はこの点に限られ、標準的な三つのaffine chart上の
消去によりschemeとしてもgraphとなる。
特に$S_q$のpoint schemeは二次元の$\mathbb P^2$である。

もし上記のpointed exact同値があれば、補題
\ref{lem:anticanonical-quadratic-hull}により$B_{\mathbf a}\simeq P$となる。
唯一の中心一次直線は保存されるので$S_q\simeq P/(t)$を得るが、
次数付き同型はpoint schemeを保つため、上の次元差に矛盾する。
\end{proof}

\begin{remark}[除外の範囲]\label{rem:aou-exclusion-scope}
この命題は、AOUの\S9で指定された無限位数のblow-up表示を除外する。
有限位数の特殊化、別の構造対象や別の標識を持つAS曲面、
任意の導来同値を除外していない。
同文献Theorem~9.4はblow-up側とAS代数側の\emph{導来同値}であり、
ここで必要なpointed exact同値と置き換えることはできない。
また第\ref{sec:integrate}節の横断方向では中心一次商のpoint schemeも
楕円曲線へ変化するため、この障害を八次元変形族全体へ延長できない。
したがって、七次元族あるいは八次元族の一般元について、
「AS曲面である」「どの標識でもAS曲面にならない」のいずれも本節からは結論しない。
公表に向けた次の比較では、許容する構造対象・標識・同値の種類と、
AOU/Okawa--Ueda側の有限位数部分の適用範囲を先に固定する必要がある。
\end{remark}
